\documentclass[10pt,a4paper]{amsart}
\usepackage[margin=19mm]{geometry}
\usepackage{amsmath,amssymb,mathtools}
\usepackage[scr=rsfs,cal=euler]{mathalfa}
\usepackage{xfrac}
\usepackage[unicode,hidelinks,bookmarks=false]{hyperref}
\newcommand{\ie}{i.e.\ }
\newcommand{\cf}{cf.\ }
\newcommand{\resp}{respectively\ }
\newcommand{\Subsection}{\subsection}
\providecommand{\nobreakdash}{\nobreak\hbox{-}}
\setlength{\emergencystretch}{2em}
\multlinegap0pt
\usepackage{bm}
\usepackage{bbm}
\usepackage{dsfont}
\usepackage{xspace}
\usepackage{cancel}
\usepackage{mathtools}
\usepackage[shortlabels]{enumitem}
\usepackage{amsthm}
\makeatletter
\@ifundefined{theorem}{\newtheorem{theorem}{Theorem}}{}
\@ifundefined{assumption}{\newtheorem{assumption}[theorem]{Assumption}}{}
\@ifundefined{corollary}{\newtheorem{corollary}[theorem]{Corollary}}{}
\@ifundefined{lemma}{\newtheorem{lemma}[theorem]{Lemma}}{}
\@ifundefined{proposition}{\newtheorem{proposition}[theorem]{Proposition}}{}
\makeatother
\DeclareMathOperator{\dist}{dist}
\DeclareMathOperator{\rank}{rank}
\newcommand{\Afib}{\mathsf{A}}
\newcommand{\C}{\mathbb{C}}
\newcommand{\R}{\mathbb{R}}
\newcommand{\cA}{\mathcal{A}}
\newcommand{\cI}{\mathcal{I}}
\newcommand{\cL}{\mathcal{L}}
\newcommand{\cT}{\mathcal{T}}
\newcommand{\dd}{\,\mathrm{d}}
\newcommand{\dvol}{\,\mathrm{dvol}}
\newcommand{\eps}{\varepsilon}
\newcommand{\fD}{\mathfrak{D}}
\newcommand{\ind}{\mathbf{1}}
\newcommand{\inner}[2]{\left\langle #1,#2\right\rangle}
\title[Laplace Asymptotics near Stratified Minimum Sets]{Laplace Asymptotics near Stratified Minimum Sets}
\author{Pavel Ievlev}
\date{Source: arXiv:2608.02626v1 (CC BY 4.0)}
\begin{document}
\maketitle

\noindent\emph{Source and licence.} Laplace Asymptotics near Stratified Minimum Sets. Version arXiv:2608.02626v1 (CC BY 4.0), distributed under the Creative Commons Attribution 4.0 International licence, as recorded in the arXiv licence metadata for this exact version and shown on the abstract page https://arxiv.org/abs/2608.02626v1. Full rights notice: licenses/arxiv-2608.02626-CC-BY-4.0.txt.\par\smallskip

\noindent\textit{Editorial scope.} Counted unit: the Proposition whose source label is \texttt{prop:robot-posterior-splitting} in the article (source0.tex lines 3786--3852, source label \texttt{prop:robot-posterior-splitting}), delivered together with the article's own complete original proof of it (source0.tex lines 3854--4187, one \texttt{proof} environment of 334 lines). No mathematics is written, repaired or re-derived: statement and proof are the article's own text line for line. Required context retained uncounted: ass:global (lines 436--447); ass:tube (lines 538--555); lem:moments (lines 759--781); eq:G-homogeneous (lines 762--764); eq:G-sphere-bounds (lines 766--769); eq:uniform-moments (lines 776--779); ass:leading-model (lines 1083--1128); cor:coercive-leading-piece (lines 1374--1394); eq:additive-leading (lines 1388--1390); thm:integrated-pointwise-cells (lines 1603--1648); eq:piece-integrated-leading (lines 1739--1741); thm:global-leading (lines 1756--1779); eq:global-additive (lines 1765--1769); eq:robot-forward-map (lines 3578--3581); eq:robot-shape-factorization (lines 3619--3625); eq:robot-minimum-set (lines 3636--3641); eq:robot-regular-arc (lines 3643--3648); lem:transverse-arc-normal-cells (lines 3670--3728); eq:transverse-arc-directional-convergence (lines 3715--3720); eq:transverse-arc-weighted-cell-limit (lines 3723--3727). Every definition, display and label of those lines is retained unchanged. Omitted from this package and not redistributed: 1--435, 448--537, 556--758, 765--765, 770--775, 780--1082, 1129--1373, 1391--1602, 1649--1738, 1742--1755, 1770--3577, 3582--3618, 3626--3635, 3642--3642, 3649--3669, 3721--3722, 3728--3785, 3853--3853, 4188--5178. Nothing inside a retained excerpt is omitted or altered: every retained body is delivered exactly as the article's lines are written, so no omission receipt of any kind accompanies this package. Citations: the retained excerpts carry no citation command at all, so no bibliography entry of the article is transcribed and no reference is redistributed. Reference adaptation: none was needed and none was made; every reference inside the delivered unit resolves inside it, so no unresolved reference or citation is produced. Printed slips kept literal: no printed word, symbol, hypothesis or number of the counted statement or of its proof was altered. Layout and class: the article's own class is \texttt{amsart} with packages including fontenc, inputenc, babel, lmodern, microtype, amsmath, amssymb, amsthm, mathtools, bm, enumitem, hyperref, geometry, graphicx; the wrapper is the lane's pinned \texttt{amsart} editorial wrapper at 10pt on A4, which restates the theorem-like environments the retained text needs and adds the article's own macro definitions that the retained text needs (\texttt{\textbackslash Afib}, \texttt{\textbackslash C}, \texttt{\textbackslash R}, \texttt{\textbackslash cA}, \texttt{\textbackslash cI}, \texttt{\textbackslash cL}, \texttt{\textbackslash cT}, \texttt{\textbackslash dd}, \texttt{\textbackslash dist}, \texttt{\textbackslash dvol}, \texttt{\textbackslash eps}, \texttt{\textbackslash fD}, \texttt{\textbackslash ind}, \texttt{\textbackslash inner}, \texttt{\textbackslash rank}), with extra wrapper packages (bm, bbm, dsfont, xspace, cancel, mathtools, enumitem). The delivered numbering is the unit's own. Assets: no retained span contains a figure environment, a graphics-inclusion command, a PDF-page inclusion, a table or a TikZ picture; no image file is copied into this package, the package declares no asset dependency and no image file is redistributed with it. Licence: version arXiv:2608.02626v1 of this article is distributed under the Creative Commons Attribution 4.0 International licence, as recorded in the arXiv licence metadata for this exact version and shown on the abstract page https://arxiv.org/abs/2608.02626v1. A full rights notice is delivered at licenses/arxiv-2608.02626-CC-BY-4.0.txt. Characters: every retained excerpt is pure ASCII, so the delivered engine, XeLaTeX with its default Latin Modern text font, reports no missing character.
\begin{assumption}[Global integrability and separation]\label{ass:global}
There is $z_*>0$ such that
\begin{equation}\label{eq:A0}
  \int_N |a(x)|e^{-z_*f(x)}\dvol_g(x)<\infty.
\end{equation}
Moreover, for every sufficiently small tubular radius $\eps>0$ there is
$\delta_\eps>0$ such that
\begin{equation}\label{eq:A1}
  f(x)\ge f_0+\delta_\eps
  \qquad\text{for }x\in N\setminus \cT_\eps(M).
\end{equation}
\end{assumption}

\begin{assumption}[Stratified adapted normal-cell change of variables]\label{ass:tube}
After decreasing $\eps>0$ if necessary, the following hold.
\begin{enumerate}[(T1),leftmargin=3em]
\item Each $F_\alpha$ is defined and $C^1$ on $U_\alpha\times B_\eps$, and
  $F_\alpha(p,0)=p$.
\item The sets $F_\alpha(\fD_\alpha(\eps))$ are pairwise disjoint up to
  $\dvol_g$-null sets and cover $\cT_\eps(M)$ up to a $\dvol_g$-null set.
\item There is a measurable Jacobian
  $J_\alpha:\fD_\alpha(\eps)\to(0,\infty)$ such that for every nonnegative
  measurable $h$,
  \begin{equation}\label{eq:cov}
    \int_{F_\alpha(\fD_\alpha(\eps))}h(x)\dvol_g(x)
    =\int_{P_\alpha}\int_{\Afib_\alpha(p)\cap B_\eps}
      h(F_\alpha(p,s))J_\alpha(p,s)\dd s\dvol_{P_\alpha}(p).
  \end{equation}
\end{enumerate}
All suprema over an empty set are understood as zero.
\end{assumption}

\begin{lemma}[Uniform exponential moments]\label{lem:moments}
Let $P$ be any parameter set and let
$G:P\times\R^d\to[0,\infty)$ be measurable.  Suppose
\begin{equation}\label{eq:G-homogeneous}
  G(p,D_u^{\bm\beta}t)=uG(p,t)
\end{equation}
for all $p,t,u$, and suppose that for some $0<g_-\le g_+<\infty$,
\begin{equation}\label{eq:G-sphere-bounds}
  g_-\le G(p,t)\le g_+
  \qquad\text{whenever }\rho_{\bm\beta}(t)=1.
\end{equation}
Then
\begin{equation}\label{eq:G-gauge-bounds}
  g_-\rho_{\bm\beta}(t)\le G(p,t)\le g_+\rho_{\bm\beta}(t)
  \qquad(p\in P,\ t\in\R^d),
\end{equation}
and for every $r\ge0$ and $\gamma>0$,
\begin{equation}\label{eq:uniform-moments}
  \sup_{p\in P}\int_{\R^d}(1+\rho_{\bm\beta}(t))^r
  e^{-\gamma G(p,t)}\dd t<\infty.
\end{equation}
The tails in~\eqref{eq:uniform-moments} converge to zero uniformly in $p$.
\end{lemma}

\begin{assumption}[Uniform coercive anisotropic model]\label{ass:leading-model}
There exist
\begin{itemize}[leftmargin=2.2em]
\item an anisotropy vector $\bm\beta\in(0,\infty)^d$ and
  $Q:=Q_{\bm\beta}$;
\item a measurable function $c:P\to(0,\infty)$ satisfying
  \begin{equation}\label{eq:c-bounds}
    0<c_-\le c(p)\le c_+<\infty;
  \end{equation}
\item a measurable function $G:P\times\R^d\to[0,\infty)$ satisfying
  \eqref{eq:G-homogeneous} and~\eqref{eq:G-sphere-bounds} uniformly in $p$;
\item a measurable remainder $\mathcal R$ such that, on $\Afib(p)\cap B_\eps$,
  \begin{equation}\label{eq:phase-leading}
    f(F(p,s))-f_0=c(p)G(p,s)(1+\mathcal R(p,s)),
  \end{equation}
  with $\mathcal R(p,0)=0$, and
  \begin{equation}\label{eq:R-uniform}
    \lim_{r\downarrow0}\sup_{p\in P}
    \sup_{s\in\Afib(p),\,\rho_{\bm\beta}(s)\le r}|\mathcal R(p,s)|=0;
  \end{equation}
\item a measurable family of limit cells
  $\Afib^\infty(p)\subset\R^d$ such that, for every $0<\delta<1$,
  \begin{equation}\label{eq:weighted-cell-limit}
  \begin{split}
   \sup_{p\in P}\int_{\R^d}
   \Big|&\ind_{D_{zc(p)}^{\bm\beta}(\Afib(p)\cap B_\eps)}(t)
        -\ind_{\Afib^\infty(p)}(t)\Big|\\
   &\hspace{4em}\times e^{-(1-\delta)G(p,t)}\dd t
   \longrightarrow0
   \qquad(z\to\infty).
  \end{split}
  \end{equation}
\end{itemize}
The radius $\eps$ is chosen small enough that $|\mathcal R(p,s)|\le\delta_0<1$
throughout the coordinate domain.  Finally, define the effective amplitude
\begin{equation}\label{eq:effective-amplitude}
  b(p,s):=a(F(p,s))J(p,s).
\end{equation}
Assume that $b$ is bounded on the truncated cell graph and that there is a
bounded measurable function $b_0:P\to\C$ such that
\begin{equation}\label{eq:amplitude-uniform}
  \lim_{r\downarrow0}\sup_{p\in P}
  \sup_{s\in\Afib(p),\,\rho_{\bm\beta}(s)\le r}
  |b(p,s)-b_0(p)|=0.
\end{equation}
\end{assumption}

\begin{corollary}[Coercive exponential leading term on one normal-cell piece]
\label{cor:coercive-leading-piece}
Under Assumptions~\ref{ass:tube} and~\ref{ass:leading-model}, define
\begin{equation}\label{eq:leading-profile}
  \cL(p):=b_0(p)\int_{\Afib^\infty(p)}e^{-G(p,t)}\dd t.
\end{equation}
Then
\begin{equation}\label{eq:uniform-leading}
  \sup_{p\in P}\left|
  e^{zf_0}(zc(p))^Q\cI(z,p)-\cL(p)
  \right|\longrightarrow0.
\end{equation}
Equivalently, there are functions $r_z:P\to\C$ with
$\sup_P|r_z|\to0$ such that
\begin{equation}\label{eq:additive-leading}
  \cI(z,p)=e^{-zf_0}z^{-Q}c(p)^{-Q}\bigl(\cL(p)+r_z(p)\bigr).
\end{equation}
In particular,~\eqref{eq:additive-leading} remains valid when $\cL(p)=0$.  If
$\inf_{p\in P}|\cL(p)|>0$, then the corresponding asymptotic equivalence is
uniform in $p$.
\end{corollary}

\begin{theorem}[Integrated leading term under pointwise cell convergence and an integrable base majorant]
\label{thm:integrated-pointwise-cells}
Fix one normal-cell piece with a common anisotropy $\bm\beta$ and put
$Q=Q_{\bm\beta}$.  Let $c:P\to(0,\infty)$ be measurable, with no uniform lower
bound assumed, and define
\[
 b(p,s):=a(F(p,s))J(p,s),\qquad
 \phi(p,s):=\frac{f(F(p,s))-f_0}{c(p)}.
\]
Let $G:P\times\R^d\to[0,\infty)$ be homogeneous and satisfy the uniform
coercivity bounds of Lemma~\ref{lem:moments}.  Suppose that there are a
measurable function $b_0:P\to\C$, measurable limit cells
$\Afib^\infty(p)\subset\R^d$, and constants $\kappa_G>0$ and
$0<\gamma\le\min\{1,\kappa_G\}$ with the following properties.
For almost every $p$,
\begin{equation}\label{eq:pointwise-cell-local-limits}
 b(p,s)\longrightarrow b_0(p),\qquad
 \frac{\phi(p,s)}{G(p,s)}\longrightarrow1
 \quad\text{as }s\to0\text{ within }\Afib(p),
\end{equation}
and throughout the truncated cells,
\begin{equation}\label{eq:pointwise-cell-lower}
 \phi(p,s)\ge\kappa_G G(p,s).
\end{equation}
Moreover, for almost every $p$,
\begin{equation}\label{eq:pointwise-weighted-cell-limit}
 \int_{\R^d}
 \left|\ind_{D_{zc(p)}^{\bm\beta}(\Afib(p)\cap B_\eps)}(t)
       -\ind_{\Afib^\infty(p)}(t)\right|
 e^{-\gamma G(p,t)}\dd t\longrightarrow0.
\end{equation}
Finally, suppose that there is a measurable $B:P\to[0,\infty)$ such that
\begin{equation}\label{eq:pointwise-cell-majorant}
 |b(p,s)|\le B(p),\qquad |b_0(p)|\le B(p),\qquad
 B(p)c(p)^{-Q}\in L^1(P,\dvol_P).
\end{equation}
Then, with
\begin{equation}\label{eq:pointwise-cell-profile}
 \cL(p):=b_0(p)\int_{\Afib^\infty(p)}e^{-G(p,t)}\dd t,
\end{equation}
one has
\begin{equation}\label{eq:pointwise-cell-result}
 I_\alpha(z)=e^{-zf_0}z^{-Q}
 \left[\int_P c(p)^{-Q}\cL(p)\dvol_P(p)+o(1)\right].
\end{equation}
\end{theorem}

\begin{equation}\label{eq:piece-integrated-leading}
 I_\alpha(z)=e^{-zf_0}z^{-Q_\alpha}\bigl(K_\alpha+o(1)\bigr).
\end{equation}

\begin{theorem}[Global leading term]\label{thm:global-leading}
Under Assumptions~\ref{ass:global} and~\ref{ass:tube}, and assuming the
piecewise expansions~\eqref{eq:piece-integrated-leading}, let
\begin{equation}\label{eq:Qstar}
  Q_*:=\min_{\alpha\in\cA}Q_\alpha,
  \qquad
  K_*:=\sum_{\alpha:Q_\alpha=Q_*}K_\alpha.
\end{equation}
Then
\begin{equation}\label{eq:global-additive}
 I(z)=e^{-zf_0}\sum_{\alpha\in\cA}
 z^{-Q_\alpha}\bigl(K_\alpha+r_\alpha(z)\bigr)
 +O\bigl(e^{-z(f_0+\delta_\eps)}\bigr),
\end{equation}
where $r_\alpha(z)\to0$ for every $\alpha$.  If $K_*\ne0$, then
\begin{equation}\label{eq:global-equivalence}
  I(z)\sim e^{-zf_0}z^{-Q_*}K_*.
\end{equation}
If $K_*=0$,~\eqref{eq:global-additive} yields only
$I(z)=o(e^{-zf_0}z^{-Q_*})$.  Identifying the first nonzero term requires
higher-order expansions on the pieces with $Q_\alpha=Q_*$; the leading
terms of pieces with larger $Q_\alpha$ cannot in general be compared with
the unspecified remainders of the dominant pieces.
\end{theorem}

\begin{equation}\label{eq:robot-forward-map}
 X(\bm q)
 =2e^{iq_1}+e^{i(q_1+q_2)}+2e^{i(q_1+q_2+q_3)}\in\C.
\end{equation}

\begin{equation}\label{eq:robot-shape-factorization}
 |S(u,v)|^2-1
 =8\cos\frac{u+v}{2}
 \left(
  \cos\frac{u-v}{2}+2\cos\frac{u+v}{2}
 \right).
\end{equation}

\begin{equation}\label{eq:robot-minimum-set}
 M:=\{\bm q\in\Omega:X(\bm q)=1\}
 =\{\bm q_0\}\sqcup\Gamma,
 \qquad
 \bm q_0:=(0,0,\pi),
\end{equation}

\begin{equation}\label{eq:robot-regular-arc}
 \Gamma:=\left\{
  \bigl(\vartheta(u),u,v_\Gamma(u)\bigr):
  \frac{4\pi}{3}\le u\le\frac{3\pi}{2}
 \right\}.
\end{equation}

\begin{lemma}[Relative normal cells for a transverse embedded arc]
\label{lem:transverse-arc-normal-cells}
Let $N$ be an $n$-dimensional Riemannian manifold, let $\Omega\subset N$ be a
relative domain, and let $\Gamma=\gamma([0,L])$ be a compact $C^2$ embedded
unit-speed arc.  Assume that $\Gamma^\circ\subset\operatorname{int}\Omega$.
For each endpoint $e$, let $\tau_e$ be the unit tangent pointing from $e$ into
$\Gamma^\circ$, and suppose that near $e$ there is a $C^2$ defining function
$h_e$ such that
\[
 \Omega=\{h_e\ge0\},\qquad \dd h_e(e)\ne0,
 \qquad \dd h_e(e)[\tau_e]>0.
\]
Then, after decreasing $\eps>0$, the nearest-point projection
$\pi_\Gamma:\cT_\eps(\Gamma)\to\Gamma$ is well defined and
$\Omega\cap\cT_\eps(\Gamma)$ has a relative normal-cell decomposition with
base pieces $\Gamma^\circ$ and the two endpoints.

More precisely, choose a $C^1$ orthonormal normal frame
$E_p:\R^{n-1}\to N_p\Gamma$ along the arc.  Over $\Gamma^\circ$ set
\[
 F_\Gamma(p,s):=\exp_p(E_p s),\qquad
 \Afib_\Gamma(p):=
 \{s\in\R^{n-1}:|s|<\eps,\ F_\Gamma(p,s)\in\Omega,\
                    \pi_\Gamma(F_\Gamma(p,s))=p\}.
\]
At an endpoint $e$, identify $T_eN$ with $\R^n$ by an orthonormal frame and set
\[
 F_e(v):=\exp_e(v),\qquad
 \Afib_e:=\{v\in T_eN:|v|<\eps,\ F_e(v)\in\Omega,\
                      \pi_\Gamma(F_e(v))=e\}.
\]
The cells are measurable.  Their three images are pairwise disjoint up to
null sets, cover $\Omega\cap\cT_\eps(\Gamma)$ up to null sets, and satisfy the
corresponding change-of-variables formulas.  Their Jacobians equal one on the
zero sections.
For each fixed $p\in\Gamma^\circ$, the scalar blow-ups of
$\Afib_\Gamma(p)$ converge locally to $\R^{n-1}$.  At an endpoint,
\begin{equation}\label{eq:transverse-arc-endpoint-cone}
 r^{-1}\Afib_e\longrightarrow
 C_e:=\{v\in T_eN:\inner{v}{\tau_e}\le0,\
                 \dd h_e(e)[v]\ge0\}
 \qquad(r\downarrow0)
\end{equation}
locally in measure.  In addition, the directions in the endpoint cell approach
those of the limiting cone:
\begin{equation}\label{eq:transverse-arc-directional-convergence}
 \sup_{\substack{v\in\Afib_e\\0<|v|\le r}}
 \dist\!\left(\frac{v}{|v|},
 C_e\cap\{w\in T_eN:|w|=1\}\right)\longrightarrow0
 \qquad(r\downarrow0).
\end{equation}
Moreover, for every $a>0$ and every continuous coercive function
$G_e:T_eN\to[0,\infty)$ satisfying $G_e(rv)=r^2G_e(v)$,
\begin{equation}\label{eq:transverse-arc-weighted-cell-limit}
 \int_{T_eN}
 \left|\ind_{r^{-1}\Afib_e}(v)-\ind_{C_e}(v)\right|
 e^{-aG_e(v)}\dd v\longrightarrow0.
\end{equation}
\end{lemma}

\begin{proposition}[A solution curve and an isolated joint-limit fold contribute
at the same order]\label{prop:robot-posterior-splitting}
For every continuous test function $\psi:\Omega\to\C$, put
\[
 Z_z[\psi]:=\int_\Omega
 \psi(\bm q)\varpi(\bm q)
 \exp\left\{-\frac z2|X(\bm q)-1|^2\right\}\dd \bm q.
\]
Thus $Z_z=Z_z[1]$.  Then
\begin{equation}\label{eq:robot-test-asymptotic}
 Z_z[\psi]=z^{-1}\left[
 2\pi\int_\Gamma
 \frac{\psi(p)\varpi(p)}{J_X(p)}\dd\ell(p)
 +C_{\mathrm{fold}}\,\psi(\bm q_0)\varpi(\bm q_0)
 +o(1)\right],
\end{equation}
where $\dd\ell$ is arclength for the product metric and
\begin{equation}\label{eq:robot-fold-constant}
 \begin{split}
 C_{\mathrm{fold}}
 &:=\int_{\R\times\R_+^2}
 \exp\left\{-\frac12\left[
 s^2+\bigl(u^2+4uv+3v^2\bigr)^2
 \right]\right\}\dd s\dd u\dd v\\
 &=\frac{\pi}{4}\log 3.
 \end{split}
\end{equation}
Consequently, if
\begin{align}
 K_{\mathrm{reg}}
 &:=2\pi\int_\Gamma
 \frac{\varpi(p)}{J_X(p)}\dd\ell(p),
 \label{eq:robot-Kreg}\\
 K_{\mathrm{fold}}
 &:=\frac{\pi}{4}\log 3\,\varpi(\bm q_0),
 \label{eq:robot-Kfold}
\end{align}
then
\begin{equation}\label{eq:robot-evidence}
 Z_z=z^{-1}\bigl(K_{\mathrm{reg}}+K_{\mathrm{fold}}+o(1)\bigr).
\end{equation}
Moreover, $\Pi_z$ converges weakly to the probability measure $\Pi_\infty$
defined by
\begin{equation}\label{eq:robot-weak-limit}
 \int_\Omega\psi\dd\Pi_\infty
 =\frac{
 \displaystyle
 2\pi\int_\Gamma
 \frac{\psi(p)\varpi(p)}{J_X(p)}\dd\ell(p)
 +C_{\mathrm{fold}}\psi(\bm q_0)\varpi(\bm q_0)}
 {K_{\mathrm{reg}}+K_{\mathrm{fold}}}.
\end{equation}
In particular,
\begin{equation}\label{eq:robot-atom-mass}
 \Pi_\infty(\{\bm q_0\})
 =\frac{K_{\mathrm{fold}}}
 {K_{\mathrm{reg}}+K_{\mathrm{fold}}}>0.
\end{equation}
For the flat prior $\varpi\equiv1$, the coefficients are explicit:
\begin{equation}\label{eq:robot-flat-prior-coefficients}
 K_{\mathrm{reg}}=\frac{\pi}{2}\log3,
 \qquad
 K_{\mathrm{fold}}=\frac{\pi}{4}\log3,
 \qquad
 \Pi_\infty(\{\bm q_0\})=\frac13.
\end{equation}
\end{proposition}

\begin{proof}
\smallskip
\noindent\emph{Step 1: Exact configurations and regularity.}
We first determine all exact configurations.  In the allowed rectangle for
the internal joint angles,
$u+v\in[\pi,17\pi/6]$.  Hence the first factor in
\eqref{eq:robot-shape-factorization} vanishes only when $u+v=\pi$, which gives
$(u,v)=(0,\pi)$.  Write
\[
 \Xi(u,v):=\cos\frac{u-v}{2}+2\cos\frac{u+v}{2}
 =3\cos\frac u2\cos\frac v2-\sin\frac u2\sin\frac v2.
\]
At $v=\pi$, the equation $\Xi=0$ again gives $u=0$.  For $v>\pi$, no solution
with $u\le\pi$ is possible.  If $u>\pi$, then $\Xi=0$ is equivalent to
\begin{equation}\label{eq:robot-tangent-equation}
 \tan\frac u2\tan\frac v2=3.
\end{equation}
On the prescribed intervals,~\eqref{eq:robot-tangent-equation} has a solution
precisely for $4\pi/3\le u\le3\pi/2$, and that solution is
$v=v_\Gamma(u)$.  This proves~\eqref{eq:robot-minimum-set} and
\eqref{eq:robot-regular-arc}.

The gradient of $\Xi$ cannot vanish on $\{\Xi=0\}$.  Indeed, if both partial
derivatives vanished, then
\[
 \sin\frac{u-v}{2}=\sin\frac{u+v}{2}=0,
\]
which is incompatible with $\Xi=0$.  Along $\Gamma$, the other factor
$\cos((u+v)/2)$ in~\eqref{eq:robot-shape-factorization} is nonzero, so
$\nabla_{u,v}|S|^2\ne0$.  At an exact solution,
$\partial_{q_1}X=i$ is tangent to the unit circle in task space, whereas a
nonzero derivative of $|S|^2$ supplies a radial task direction.  Thus
$\rank DX=2$ on $\Gamma$, and $J_X$ is bounded away from zero
there.  At $(u,v)=(4\pi/3,4\pi/3)$, implicit differentiation gives
$v_\Gamma'(u)=-1$; the other endpoint is parametrized with nonzero $u$-velocity at
$u=3\pi/2$.  Hence the two intersections with the joint-limit faces are
transverse.

\smallskip
\noindent\emph{Step 2: Localization near the two solution components.}
Choose disjoint neighborhoods $U_\Gamma$ and $U_0$ of $\Gamma$ and $\bm q_0$,
with smooth cutoffs $\chi_\Gamma$ and $\chi_0$ equal to one near the respective
components.  Since $\Omega$ is compact and~\eqref{eq:robot-minimum-set} lists
all exact solutions, the remaining part of $Z_z[\psi]$ is exponentially
small.

\smallskip
\noindent\emph{Step 3: The regular solution curve.}
For the regular component, put
\[
 e_-:=\bigl(\vartheta(4\pi/3),4\pi/3,4\pi/3\bigr),
 \qquad
 e_+:=\bigl(\vartheta(3\pi/2),3\pi/2,v_\Gamma(3\pi/2)\bigr),
 \qquad
 \Gamma^\circ:=\Gamma\setminus\{e_-,e_+\}.
\]
We verify the hypotheses of the normal-cell theorems rather than derive this
contribution by a separate coarea asymptotic.  The transversality established in
Step~1 permits the active joint-limit defining functions to be oriented so that
$\dd h_e(e)[\tau_e]>0$.  Lemma~\ref{lem:transverse-arc-normal-cells} then supplies
ordinary nearest-point normal cells over $\Gamma^\circ$ and two endpoint caps,
assigned to the zero-dimensional pieces $\{e_-\}$ and $\{e_+\}$.  These pieces
are almost everywhere disjoint, cover a sufficiently small relative
neighborhood of $\Gamma$, and satisfy the relative-domain form of
Assumption~\ref{ass:tube}.

Choose a $C^1$ orthonormal normal frame along $\Gamma^\circ$ and write
$E_p:\R^2\to N_p\Gamma$ for the associated isometry.  Define
\begin{equation}\label{eq:robot-regular-Hessian}
 H_\Gamma(p):=E_p^{\mathsf T}DX(p)^{\mathsf T}DX(p)E_p,
 \qquad
 G_\Gamma(p,s):=\frac12\inner{H_\Gamma(p)s}{s}.
\end{equation}
Since $\ker DX(p)=T_p\Gamma$ and $DX(p)$ has rank two on the compact arc,
$H_\Gamma(p)$ is uniformly positive definite up to both endpoints.  Taylor's
theorem in normal coordinates gives, uniformly as $s\to0$,
\begin{equation}\label{eq:robot-regular-normal-model}
 \frac12|X(F(p,s))-1|^2
 =G_\Gamma(p,s)+O(|s|^3),
 \qquad p\in\Gamma^\circ.
\end{equation}
For each fixed $p\in\Gamma^\circ$, the point lies in the interior of $\Omega$,
so its relative normal cell contains a neighborhood of zero in $\R^2$.
Consequently, under the scalar quadratic dilation,
\[
 z^{1/2}\Afib_\Gamma(p)\longrightarrow\R^2
\]
in the weighted-indicator sense of
Theorem~\ref{thm:integrated-pointwise-cells}.  This convergence need not be
uniform as $p$ approaches an endpoint, which is precisely why the pointwise
cell theorem is the appropriate result here.

Set $A_\Gamma=\chi_\Gamma\psi\varpi$.  In the corresponding Fermi
coordinates,
\[
 A_\Gamma(F(p,s))J(p,s)\longrightarrow\psi(p)\varpi(p).
\]
The amplitude is bounded, the base $\Gamma^\circ$ has finite length, and the
uniform ellipticity of~\eqref{eq:robot-regular-Hessian} supplies both the
coercive lower bound and an integrable Gaussian majorant.  Thus all hypotheses
of Theorem~\ref{thm:integrated-pointwise-cells} hold with
$\bm\beta=(1/2,1/2)$, $Q=1$, $c\equiv1$, and limit cell $\R^2$.  Denoting the
contribution of this normal-cell piece by $I_{\Gamma^\circ}(z)$,
Theorem~\ref{thm:integrated-pointwise-cells} gives
\begin{equation}\label{eq:robot-regular-interior-asymptotic}
 I_{\Gamma^\circ}(z)
 =z^{-1}\left[
 \int_{\Gamma}
 \psi(p)\varpi(p)
 \int_{\R^2}e^{-G_\Gamma(p,t)}\dd t\,\dd\ell(p)
 +o(1)\right],
\end{equation}
where adding the two endpoints to the outer integral does not change it.
Because $DX(p)$ annihilates the unit tangent to $\Gamma$, an orthonormal
splitting into that tangent and the columns of $E_p$ gives
\[
 \det H_\Gamma(p)
 =\det\bigl(DX(p)DX(p)^{\mathsf T}\bigr)
 =J_X(p)^2.
\]
The Gaussian integral in~\eqref{eq:robot-regular-interior-asymptotic} is
therefore $2\pi/J_X(p)$.

\smallskip
\noindent\emph{Step 4: Endpoint caps of the regular curve.}
It remains only to check that the two cap pieces are lower order.  Fix an
endpoint $e\in\{e_-,e_+\}$, let $\tau_e$ be the unit tangent to $\Gamma$
pointing from $e$ into $\Gamma^\circ$, and let $h_e$ be a local defining
function for the active joint-limit face, oriented so that
$\Omega=\{h_e\ge0\}$.  Transversality gives
$\dd h_e(e)[\tau_e]>0$.  Lemma~\ref{lem:transverse-arc-normal-cells} identifies the scalar blow-up of
the endpoint normal cell as
\begin{equation}\label{eq:robot-endpoint-cone}
 C_e:=\{v\in T_e\mathsf Q:
       \inner{v}{\tau_e}\le0,\ \dd h_e(e)[v]\ge0\}.
\end{equation}
The first inequality is the nearest-point end-cap condition and the second is
feasibility.  Since $\ker DX(e)=\R\tau_e$, the two inequalities and
$\dd h_e(e)[\tau_e]>0$ imply
\[
 C_e\cap\ker DX(e)=\{0\}.
\]
Hence $|DX(e)v|^2$ is bounded below by a positive multiple of $|v|^2$ on
$C_e$.  Choose a closed conic neighborhood $\mathcal V_e$ of $C_e\setminus\{0\}$ that
is still disjoint from the unit kernel directions.  By
\eqref{eq:transverse-arc-directional-convergence}, all nonzero directions in a
sufficiently small endpoint cell lie in $\mathcal V_e$.  We use the homogeneous model
\[
 \widetilde G_e(v)
 :=\frac12|DX(e)v|^2+\dist(v,\mathcal V_e)^2
\]
on the ambient tangent space.  It is coercive everywhere, and on the actual
cell the distance term vanishes, so $\widetilde G_e(v)=|DX(e)v|^2/2$ there.
Taylor expansion gives
\[
 \frac12|X(\exp_e v)-1|^2
 =\widetilde G_e(v)\bigl(1+O(|v|)\bigr)
\]
within that cell, while
\eqref{eq:transverse-arc-weighted-cell-limit} supplies the required weighted
convergence of the scaled cell to $C_e$.  The relative-domain form of
Corollary~\ref{cor:coercive-leading-piece}, with model
$\widetilde G_e$, three quadratic weights, and $Q=3/2$, therefore gives
\[
 I_e(z)=O(z^{-3/2}).
\]
Applying this to both endpoints and combining it with
\eqref{eq:robot-regular-interior-asymptotic} gives the regular contribution
\begin{equation}\label{eq:robot-regular-asymptotic}
 \int_\Omega A_\Gamma(\bm q)e^{-z|X(\bm q)-1|^2/2}\dd \bm q
 =z^{-1}\left[
 2\pi\int_\Gamma\frac{\psi(p)\varpi(p)}{J_X(p)}\dd\ell(p)
 +o(1)\right].
\end{equation}
This piece has exponent $Q_\Gamma=1$; its two endpoint strata have the larger
exponent $3/2$ and do not enter the leading coefficient.

\smallskip
\noindent\emph{Step 5: The isolated joint-limit fold.}
It remains to analyze the isolated exact configuration at the joint-limit
corner.  Near $\bm q_0$, use the following linear coordinates.  Their determinant
is one, so they preserve the volume element:
\begin{equation}\label{eq:robot-corner-coordinates}
 u:=q_2,
 \qquad
 v:=q_3-\pi,
 \qquad
 s:=q_1-u-2v.
\end{equation}
After shrinking the coordinate neighborhood, the relative cell is exactly
\[
 C_0:=\R\times\R_+^2
\]
up to the Euclidean truncation $B_\eps$; the upper joint limits are inactive
there.  The adapted map
\[
 F_0(s,u,v)=(s+u+2v,u,\pi+v)
\]
is linear with volume Jacobian one.  Assumption~\ref{ass:tube} does not require
an orthonormal normal frame, so the normal-cell framework applies directly in
these shear coordinates.  In particular, no auxiliary orthonormalization or
moving-cell comparison is needed at the corner.

Put
\begin{equation}\label{eq:robot-corner-quadratic}
 \mathcal P_{\mathrm{fold}}(u,v):=u^2+4uv+3v^2=(u+v)(u+3v).
\end{equation}
For $h\downarrow0$, Taylor expansion of~\eqref{eq:robot-forward-map} gives,
uniformly for $(\xi,U,V)$ in compact subsets of
$\R\times\R_+^2$,
\begin{equation}\label{eq:robot-weighted-residual}
 \begin{split}
 &X\bigl(h^2\xi+hU+2hV,hU,\pi+hV\bigr)-1\\
 &\hspace{7em}
 =h^2\bigl(\mathcal P_{\mathrm{fold}}(U,V)+i\xi\bigr)+O(h^3).
 \end{split}
\end{equation}
Equivalently, with the weighted quasi-radius
\[
 r_{\mathrm{fold}}(s,u,v):=\bigl[s^2+(u^2+v^2)^2\bigr]^{1/4},
\]
the unscaled residual admits the uniform estimate
\begin{equation}\label{eq:robot-weighted-residual-bound}
 X(s+u+2v,u,\pi+v)-1=\mathcal P_{\mathrm{fold}}(u,v)+is+\mathcal R_{\mathrm{fold}}(s,u,v),
 \qquad |\mathcal R_{\mathrm{fold}}(s,u,v)|\le C r_{\mathrm{fold}}(s,u,v)^3
\end{equation}
on the feasible cone near the origin.  This follows either by grouping the
ordinary Taylor series by the weights $(1/2,1/4,1/4)$ or by applying
\eqref{eq:robot-weighted-residual} on the compact weighted unit sphere.
Thus the squared-residual phase has weighted principal model
\begin{equation}\label{eq:robot-principal-model}
 G(s,u,v):=\frac12\left[s^2+\mathcal P_{\mathrm{fold}}(u,v)^2\right],
 \qquad
 \bm\beta=\left(\frac12,\frac14,\frac14\right),
 \qquad Q=1.
\end{equation}
On $\R\times\R_+^2$,
$\mathcal P_{\mathrm{fold}}(u,v)\ge u^2+3v^2$, so $G$ is comparable to $r_{\mathrm{fold}}(s,u,v)^4$ and is coercive
on the admissible cone.  From~\eqref{eq:robot-weighted-residual-bound},
$|\mathcal P_{\mathrm{fold}}(u,v)+is|=O(r_{\mathrm{fold}}^2)$ and $|\mathcal R_{\mathrm{fold}}|=O(r_{\mathrm{fold}}^3)$; hence the cross term in the squared
residual is $O(r_{\mathrm{fold}}^5)$ and $|\mathcal R_{\mathrm{fold}}|^2=O(r_{\mathrm{fold}}^6)$.  Dividing by the lower bound
$G\ge cr_{\mathrm{fold}}^4$ gives the uniform relative estimate
\begin{equation}\label{eq:robot-phase-relative-error}
 \frac12|X(\bm q)-1|^2
 =G(s,u,v)\left(1+O\left(
 [s^2+(u^2+v^2)^2]^{1/4}
 \right)\right)
\end{equation}
on the cone near the origin.  Although $\mathcal P_{\mathrm{fold}}$ is not coercive on all of $\R^2$,
Corollary~\ref{cor:coercive-leading-piece} requires only an ambient homogeneous
extension agreeing with the phase model on $C_0$.  Put
\begin{equation}\label{eq:robot-extended-model}
 \widehat G(s,u,v):=
 \frac12\left[
 s^2+\mathcal P_{\mathrm{fold}}(u_+,v_+)^2+u_-^4+v_-^4
 \right],
\end{equation}
with the standing positive- and negative-part notation.  This function is continuous,
nonnegative, coercive, and homogeneous for the weights in
\eqref{eq:robot-principal-model}.  On the admissible cell $C_0$ it coincides
with $G$.  The cell $C_0$ is anisotropically invariant, while the truncation by
$B_\eps$ recedes to infinity after scaling.  The relative-domain form of
Corollary~\ref{cor:coercive-leading-piece}, applied in the determinant-one
coordinates
\eqref{eq:robot-corner-coordinates}, therefore gives
\begin{equation}\label{eq:robot-fold-asymptotic}
 \int_\Omega\chi_0(\bm q)\psi(\bm q)\varpi(\bm q)
 e^{-z|X(\bm q)-1|^2/2}\dd \bm q
 =z^{-1}\left[
 C_{\mathrm{fold}}\psi(\bm q_0)\varpi(\bm q_0)+o(1)
 \right].
\end{equation}
This is the step for which the anisotropic scaling is essential.  At
$\bm q_0$, the ordinary Hessian of the squared error has rank one: only the $s$
direction is visible to second order.  The two inward joint motions $u,v$
first appear through the quartic term $\mathcal P_{\mathrm{fold}}(u,v)^2$, so an ordinary
three-dimensional Gaussian Laplace approximation cannot produce the correct
coefficient.

\smallskip
\noindent\emph{Step 6: Coefficient evaluation and global comparison.}
It remains to evaluate the model integral.  Integrating first in $s$ and then
putting $v=tu$ in the first quadrant gives
\begin{align*}
 C_{\mathrm{fold}}
 &=\sqrt{2\pi}\int_{\R_+^2}
 e^{-\mathcal P_{\mathrm{fold}}(u,v)^2/2}\dd u\dd v\\
 &=\sqrt{2\pi}\,\frac{\sqrt{2\pi}}4
 \int_0^\infty\frac{\dd t}{(1+t)(1+3t)}\\
 &=\frac{\pi}{4}\log3,
\end{align*}
which proves~\eqref{eq:robot-fold-constant}.  The regular interior piece
and the isolated fold both have exponent $Q=1$, whereas the two endpoint caps
have exponent $3/2$ and the complement is exponentially small.  The global
comparison theorem, Theorem~\ref{thm:global-leading}, therefore sums
\eqref{eq:robot-regular-asymptotic} and
\eqref{eq:robot-fold-asymptotic} and proves
\eqref{eq:robot-test-asymptotic}.  Taking $\psi=1$ gives
\eqref{eq:robot-evidence}, and division of the two expansions proves
\eqref{eq:robot-weak-limit} and~\eqref{eq:robot-atom-mass}.

\smallskip
\noindent\emph{Step 7: Flat-prior constants.}
For the flat prior, the regular coefficient supplied by
Theorem~\ref{thm:integrated-pointwise-cells} can be evaluated explicitly.
Put $\Delta(u,v):=|S(u,v)|^2-1$.  At this stage coarea is used
only to compute that already established geometric coefficient, not to derive
a separate asymptotic formula.  Since $q_1$ only rotates the shape vector,
polar task coordinates and the coarea formula give the following identity,
where $\dd\ell_{uv}$ denotes Euclidean arclength in the $(u,v)$-plane:
\begin{equation}\label{eq:robot-flat-regular-reduction}
 \int_\Gamma\frac{\dd\ell}{J_X}
 =\int_{\{\Delta=0\}\cap\mathrm{regular\ shape\ arc}}
 \frac{2}{|\nabla\Delta|}\dd\ell_{uv}
 =2\int_{4\pi/3}^{3\pi/2}
 \frac{\dd u}{|\partial_v\Delta(u,v_\Gamma(u))|}.
\end{equation}
Set $x=-\tan(u/2)$.  Along the arc,
$-\tan(v_\Gamma(u)/2)=3/x$, and direct differentiation of
\eqref{eq:robot-shape-factorization} gives
\[
 |\partial_v\Delta(u,v_\Gamma(u))|=\frac{8x}{1+x^2},
 \qquad
 \dd u=-\frac{2\dd x}{1+x^2}.
\]
As $u$ increases from $4\pi/3$ to $3\pi/2$, $x$ decreases from $\sqrt3$ to
$1$.  Hence
\[
 \int_\Gamma\frac{\dd\ell}{J_X}
 =\frac14\log3.
\]
Together with~\eqref{eq:robot-fold-constant}, this proves
\eqref{eq:robot-flat-prior-coefficients}.
\end{proof}

\end{document}
